\unit{강연 236 — 피카르 스킴: 일반 성질}
\sid{236}
\src{221}
\seminar{제14년차, 1961/62, 제236호\\1962년 5월}
\paperhead{대수기하학에서의 하강 기법과 존재 정리\\
VI. 피카르 스킴: 일반 성질}%
  {알렉상드르 그로텐디크 지음}

\fsection{0}{앞 강연의 보충 ([3], V)\textsuperscript{(1)}}

[3], V에서 제기한 피카르 준스킴의 존재 문제와 관련하여 약간의 진전이
있었다.

\medskip
\noindent a. (멈퍼드). $f:X\longrightarrow S$가 사영이고 분리 가능한
사상(= 평탄하고 섬유들이 분리 가능)일 때 피카르 준스킴
$\operatorname{Pic}_{X/S}$가 존재한다는 명제는 일반적으로 참이 아니다.
이는 $f$의 섬유들이 1차원이고 $S$가 완비 이산 값매김환의 스펙트럼인
경우에도 마찬가지이다. 예를 들려면
$S=\operatorname{Spec}\mathbf R[[t]]$로 놓고, $X$를 다음 방정식으로
정의되는 $\mathbf P^2_S$의 부분스킴으로 잡으면 된다(여기서 $x,y,z$는
동차변수이다).
\[
x^2+y^2=tz^2.
\]
따라서 이 식은 기하적으로는 서로 만나는 두 직선으로 퇴화하는 이차곡선을
나타내지만, 체 $\mathbf R$ 위의 특수 섬유는 여전히 기약이다(그 섬유는
$\mathbf R$ 위의 방정식 $x^2+y^2=0$으로 주어진다). 에탈 확대
$S'\longrightarrow S$, $S'=\operatorname{Spec}\mathbf C[[t]]$를 한 뒤에는
$X'/S'$의 피카르 준스킴이 존재함을 쉽게 알 수 있으며, 이를
$\widetilde S'$의 복사본들의 서로소 합으로 명시적으로 기술할 수 있다.
여기서 $\widetilde S'$은 $S'$의 원점을 무한히 여러 번 둘로 갈라서 얻는다.
$S'\longrightarrow S$에 대한 $\operatorname{Pic}_{X'/S'}$ 위의 하강 자료는
(여기서는 $S'$의 $S$ 위 갈루아 군 $\mathbf Z/2\mathbf Z$의 작용으로
주어진다) 유효하지 않음을 쉽게 확인할 수 있다. 이 군이 둘로 갈라진 몇몇
점을 서로 교환하므로 어떤 궤도들은 아핀 열린집합에 들어 있지 않기
때문이다. 그러나 멈퍼드는 $f:X\longrightarrow S$가 사영이고 분리 가능한
사상이며, 모든 $s\in S$에 대하여 $X_s$의 기약성분들이 $k(s)$에 관하여
기하적으로 기약이면 $\operatorname{Pic}_{X/S}$가 존재함을 보일 수 있다.
그 증명은 자신의 몫 정리를 정밀화한 결과에 의존한다([7] 참조). 한편
섬유 $X_s$의 기약성분들에 아무 가정도 두지 않더라도, 아래에서 도입할
스킴 $\operatorname{Pic}^{\tau}_{X/S}$가 존재할 가능성은 여전히 남아
있음에 유의하라.

\src{222}
\noindent b. (뮈르). $X$가 한 체 위의 고유 스킴이면
$\operatorname{Pic}_{X/k}$가 존재한다. 그 증명은 슈발레의 증명 [1]을
일부 되풀이하며, 피카르 함자의 군 구조를 본질적으로 이용한다.

피카르 스킴 이론에 관한 몇 가지 보충, 특히 아벨 스킴과 관련된 내용은
[7]을 참조하는 것이 유용할 것이다. 끝으로 이 강연에는 주목할 만한
공백이 하나 있다. 사영 스킴의 피카르 스킴과 그 초평면 단면들의 피카르
스킴을 비교하게 해 주는 ``동치 판정법''이 없다는 점이다. 그러한 판정법을
전개하는 데 필요한 핵심 정리들은 [5]에 있으며, 여기에 피카르 스킴의
존재 정리들을 함께 사용해야 한다.

\fsection{1}{가환 군 준스킴의 위상적 성질}

먼저 $k$를 체, $G$를 $k$ 위의 군 준스킴이라 하자. 항등원 $e$는
$k$-유리이므로 반드시 닫혀 있고, 따라서 $G\times_kG$의 대각선은 곧바로
닫혀 있다. 그러므로 $G$는 분리되어 있다. 즉 체 위의 모든 군 준스킴은
분리되어 있다. 항등원 $e$의 연결성분을 $G^\circ$로 나타내자. $e$가
$k$-유리이므로 $G^\circ$는 실제로 기하적으로 연결되어 있다. 다시 말해
$G^\circ$를 만드는 과정은 기초체의 확대와 양립한다. 또한 $G^\circ$는
집합으로서 곱셈에 안정하다. $G$가 국소 뇌터이면 $G^\circ$가 열려
있으므로, $G^\circ$를 $G$의 열린 부분군으로 볼 수 있다. 이하에서는
$G$가 $k$ 위에서 국소 유한형이라고 가정한다. 그러면 $G^\circ$는
기하적으로 기약이고 $k$ 위에서 유한형이다. 실제로 $k$가 대수적으로
닫혀 있다고 가정해도 되고, 더 나아가 $G$가 축약이라고 가정해도 된다.
(이 경우 $G_{\mathrm{réd}}\times_kG_{\mathrm{réd}}$도 여전히 축약이므로
$G_{\mathrm{réd}}$는 $G$의 부분군이다.) 따라서 $G$는 비어 있지 않은
열린집합 위에서 $k$에 대하여 매끄럽고, 그 열린집합을 평행이동하면 모든
곳에서 매끄럽다. 이제 $G$는 국소적으로 기약이므로 그 기약성분들은
연결성분들과 일치하며, 따라서 $G^\circ$는 기약이다. $U$를 $G^\circ$에서
$e$의 아핀 근방이라 하자. $G^\circ$가 기약이라는 사실을 이용하면
$U\cdot U=G^\circ$임을 곧바로 알 수 있다. 이는 $G^\circ$가
준콤팩트이고, 따라서 $k$ 위에서 유한형임을 보인다.

간단히 하기 위해 $G$가 가환이라고 가정하자. 모든 정수 $n>0$에 대하여
$G$의 $n$제곱 준동형 $\varphi_n$에 의한 $G^\circ$의 역상을
$G^{(n)}$이라 하자. 그러면 $G^{(n)}$은 $G$의 열린 부분군이다. 다음과
같이 놓는다.
\[
G^\tau=\bigcup_nG^{(n)}.
\]

\src{223}
\[
G^\sigma=\bigcup_{(n,p)=1}G^{(n)},
\qquad
G^\rho=\bigcup_hG^{(p^h)},
\]
여기서 $p$는 체 $k$의 표수 지수이다. 이로써 $G$의 열린 부분군들을
얻으며, 이들은
\[
G^\sigma\cap G^\rho=G^\circ,
\qquad
G^\sigma G^\rho=G^\tau
\]
를 만족한다.

\result{비고}{} 몫 군 스킴 $G/G^\circ=N$을 구성할 수 있다([3], IV 참조).
그러면 $G^\tau$, $G^\sigma$, $G^\rho$를 각각 $N$의 유한 위수 원소들로
이루어진 부분군, 그 $p$-일차 성분의 자연스러운 여부분(즉 소수
$q\neq p$에 대한 $q$-일차 성분들의 합), 그리고 그 $p$-일차 성분의
$G$ 안에서의 역상으로 정의할 수 있다. 이와 관련하여 $N$은 $k$ 위의
이산 분리 가능 군 스킴임에 유의하라. 따라서 $k$의 대수적 폐포
$\overline k$를 하나 고르면(이에 따라 갈루아 군 $\pi$가 생긴다), $N$은
$\pi$가 자기동형사상으로 작용하는 보통의 이산군과 동일시된다. 이
관점에서 유한 위수 원소들의 부분군을 구성하고 이를 $q$-일차 성분들로
분해하는 것을 명백하게 해석할 수 있다. $G$가 $k$ 위의 고유 스킴 $X$의
피카르 스킴일 때, $N$을 $k$ 위의 $X$의 (축약)
\emph{네론--세베리 스킴}이라 부를 수 있다. $G^\circ_{\mathrm{réd}}$가
$G^\circ$의 군 부분스킴일 때에도 몫
$N'=G/G^\circ_{\mathrm{réd}}$를 도입하는 것이 적절하다. 이는 특히
$k$가 완전체이거나 $G^\circ$가 $k$ 위에서 고유일 때(예를 들어 $X$가
기하적으로 정규일 때) 성립한다. 이 몫은 특수화의 관점, 즉 $X$가
대수적 스킴들의 족 안에서 변할 때 $N$보다 더 잘 거동하는 경향이 있다.

이제 $S$를 국소 뇌터 준스킴, $G$를 $S$ 위에서 국소 유한형인 $S$ 위의
군 준스킴이라 하자. $G$가 $S$ 위에서 유한형이거나 분리되어 있다고
가정하지 않는다. 이때
\[
G^\circ=\bigcup_{s\in S}(G_s)^\circ
\]
로 놓고, $G$가 가환일 때에는
\[
G^\tau=\bigcup_{s\in S}(G_s)^\tau,
\qquad
G^\sigma=\bigcup_{s\in S}(G_s)^\sigma,
\qquad
G^\rho=\bigcup_{s\in S}(G_s)^\rho
\]
로 놓는다.

\src{224}
이들은 $G$의 곱셈에 안정한 부분집합들이지만, 그렇다고 이들을 $G$의 군
부분준스킴으로 정의할 수 있는 것은 물론 아니다. 특히 바탕집합이
$G^\circ$인 $G$의 군 부분준스킴은 일반적으로 존재하지 않는 듯하다.
물론 이 집합들 가운데 하나가 열려 있다면, 유도된 구조를 부여했을 때
$G$의 열린 군 부분준스킴이 된다. 우리는 $G^\tau$의 경우에는 언제나
그러함을 보일 것이다. 따라서 표현 가능 함자의 관점, 특히 ``특수화''의
관점에서는 수치적 동치가 대수적 동치보다 더 만족스럽게 거동한다. 이제
막 정의한 집합들의 주된 일반 성질을 서술하자.

\result{정리}{1.1} $G^\circ$, $G^\tau$, $G^\sigma$, $G^\rho$는 국소
구성가능이다. 더 나아가 다음이 성립한다.

\noindent (i) $G^\circ$는 $S$ 위에서 준콤팩트이다. $G_s^\circ$들이
고유이고 $G$가 $S$ 위에서 분리되어 있으면 $G^\circ$는 $S$ 위에서
고유이며, 따라서 $G$에서 닫혀 있다.

\noindent (ii) $G^\tau$는 열려 있다. $G^\circ$가 닫혀 있으면
$G^\tau$도 닫혀 있다.

\noindent (iii) $G^\circ$가 닫혀 있으면, 동표수인 경우, 즉 $S$의 모든
잉여체가 같은 표수를 갖는 경우에는 $G^\sigma$도 닫혀 있다.
$G^\circ$가 닫혀 있고 $G\longrightarrow S$가 $G^\sigma$의 점들에서
보편적으로 열려 있으면(아래의 따름정리 1.5 참조),
$G^\sigma\longrightarrow S$는 보편적으로 열린다.

\noindent (iv) $G^\circ$가 닫혀 있으면 $G^\rho$도 닫혀 있다.
동표수라고 가정하고, $(n,p)=1$인 모든 정수 $n>0$에 대하여 $G$의
$n$제곱 준동형이 열린 사상이라고 가정하자. 그러면 $G^\rho$는 열려 있다.

증명의 개요를 제시하자. $G^\circ$가 국소 구성가능이고 $S$ 위에서
준콤팩트라는 사실은 다음 보조정리에 포함되어 있다.

\result{보조정리}{1.2} $S\neq\varnothing$이라고 가정하자. 그러면
$S$ 안에 비어 있지 않은 열린집합 $U$가 있고, $U$ 위에 연결인 섬유들을
갖는 유한형 군 스킴 $H$가 있으며, 상이 $G^\circ|U$인 열린 몰입 사상인
군 스킴 준동형 $H\longrightarrow G|U$가 있다.

보조정리를 증명하기 위해 $S$가 기약이라고 가정해도 된다. 그 일반점을
$\eta$라 하고 밑변경
$S'=\operatorname{Spec}(\mathcal O_{S,\eta})\longrightarrow S$를 하자.
그러면 국소 아르틴 환 $\mathcal O_{S,\eta}=A$ 위의 군 스킴 $G'$을 얻고,
앞서 말한 바와 같이 그 안에는 $A$ 위에서 유한형인 열린 군 부분스킴
$G'{}^\circ$가 있다. 따라서 이는 $\eta$의 어떤 열린 근방 $U$ 위의
유한형 군 스킴 $H$에서 유도되며, 표준 몰입 사상
$G'{}^\circ\longrightarrow G'$은 몰입 사상

\src{225}
$H\longrightarrow G|U$에서 유도된다. $U$를 충분히 작게 잡으면 이
몰입 사상은 열려 있고 군 스킴 준동형이다. 또한 $U$를 충분히 작게 잡으면
$H$의 섬유들은 연결되어 있고 모두 같은 차원을 가지며, 그 차원은 $G$의
섬유들의 차원과 같다. 따라서 모든 $s\in U$에 대하여 $G_s$ 안에서
$H_s$의 상은 정확히 $G_s^\circ$이다(필요하면 $U$를 다시 충분히 작게
잡는다). 이로써 1.2가 증명된다.

(i)의 두 번째 명제는 다음 보조정리에 포함되어 있다(이를 $G$ 안에서
$G^\circ$의 준콤팩트 열린 근방에 적용한다).

\result{보조정리}{1.3} $S$가 국소 뇌터이고, $X$가 $S$ 위에서 유한형이고
분리된 준스킴이며, $g$가 $S$ 위의 $X$의 단면이라고 하자. $X^\circ$를
각 $X_s$에서 $g(s)$를 포함하는 연결성분들의 합집합이라 하자. 어떤
$s\in S$에 대하여 $X_s^\circ$가 $k(s)$ 위에서 고유이면, $s$의 열린
근방 $U$가 있어서 $X^\circ|U$는 $U$ 위에서 고유이고, 특히 $X|U$에서
닫혀 있다.

밑공간의 충실 평탄 하강에 의해 $S$가 완비 국소환의 스펙트럼이고 $s$가
그 닫힌점인 경우로 환원된다. [2], III, 5.5.1을 적용하면 $X$가 서로소인
두 열린집합 $X'$과 $X''$의 합으로 분해됨을 알 수 있다. 여기서 첫째 것은
$S$ 위에서 고유이고 $X'_s=X_s^\circ$를 만족한다. 그러므로 $X=X'$, 즉
$X$가 $S$ 위에서 고유인 경우로 환원된다. 이 경우에는 부분집합의 고유성에
대한 값매김 판정법([2] 제II장에서 빠져 있다)을 이용하는 표준적인 증명으로
결론을 얻는다.

$G^\tau$가 열려 있음을 증명하자. $G^\tau$를 만드는 과정($G^\circ$,
$G^\sigma$, $G^\rho$의 경우도 마찬가지이다)이 밑변경과 양립함을 감안하면,
이는 모든 $S$ 위의 $G$의 단면 $g$에 대하여 $g^{-1}(G^\tau)$가 열려
있음을 보이는 것과 같다. 이는 다음 두 명제를 뜻한다.

\medskip
\noindent a. $y\in S$이고 $g(y)\in G_y^\tau$라 하자. 그러면
$y$에 충분히 가까운 모든 $y'\in\overline{\{y\}}$에 대하여
$g(y')\in G_{y'}^\tau$이다.

\medskip
\noindent b. $y'\in S$이고 $g(y')\in G_{y'}^\tau$라 하자. 그러면
$y'$의 모든 일반화 $y$에 대하여 $g(y)\in G_y^\tau$이다.

(a)를 위해 $g^n(y)\in G_y^\circ$가 되는 정수 $n>0$이 존재함에 유의하자.
$G^\circ$가 구성가능이므로 $(g^n)^{-1}(G^\circ)$도 구성가능하다. 따라서
$y$에 충분히 가까운 $y'\in\overline{\{y\}}$에 대하여
$g^n(y')\in G_{y'}^\circ$이다. (b)를 위해서는 $g(y')^n=g^n(y')$들이
$G_{y'}$의 한 준콤팩트 열린집합 안에 머묾에 유의하자. 실제로 이들은
$G_{y'}^\circ$를 법으로 한 유한 개의 잉여류 안에 들어 있다. 그러므로
$g^n(y')$들을 포함하는 $G$의 준콤팩트 열린집합 $U$가 있고, 이들의
일반화 $g^n(y)$들도 $U$에 들어 있다. 따라서 $g(y)$의 거듭제곱들은
$G_y$의 한 준콤팩트 열린집합 안에 머물며, 이로부터
$g(y)\in G_y^\tau$임을 쉽게 얻는다.

\src{226}
$G^\circ$가 닫혀 있다고 가정하고 $G^\tau$도 닫혀 있음을 증명하자.
$G^\tau$가 열려 있고 따라서 국소 구성가능임은 이미 알고 있으므로,
특수화에 안정함만 보이면 된다. 이는 $G^\tau$가 닫힌집합들의 합집합,
곧 $n$제곱 준동형 $\varphi_n$에 의한 닫힌집합 $G^\circ$의 역상들의
합집합이라는 사실에서 따른다.

$G^\circ$가 닫혀 있으면 같은 논증으로 $G^\sigma$와 $G^\rho$도 닫혀
있음을 보일 수 있다(첫째 경우에는 동표수라는 조건을 둔다). 그러려면
$G^\sigma$와 $G^\rho$가 국소 구성가능임을 먼저 보이면 된다.
$x\in G^\tau$에 대하여, $n$제곱 준동형 $\varphi_n$이 $x$를
$G^\circ$로 보내는 가장 작은 정수 $n>0$을 $\nu(x)$라 하자. 그러면
$G^\sigma$(각각 $G^\rho$)는 $\nu(x)$가 $p$와 서로소인(각각 $p$의
거듭제곱인) $x\in G$\footnote{인쇄본에는 여기에 ``$X$''라고 되어 있으나,
이 대목에서 $X$는 정의되어 있지 않다. 지금 살피는 네 장소는 모두
$G$의 부분집합이다.}로 이루어진다. 이제 구성가능성에 관한 명제는 다음의
더 정밀한 명제에서 따른다.

\result{보조정리}{1.4} $G^\tau$ 위의 함수 $\nu$는 국소 구성가능이다.

실제로 이는 모든 정수 $n>0$에 대하여 $\nu(x)=n$인 $x\in G^\tau$들의
집합이 국소 구성가능이라는 뜻이다. 그런데 이 집합은
$\varphi_n^{-1}(G^\circ)$와, $d$가 $n$의 진약수들을 움직일 때의
$\varphi_d^{-1}(G^\circ)$들의 합집합의 차집합이다. $G^\circ$가 국소
구성가능이므로 $\varphi_d^{-1}(G^\circ)$들도 국소 구성가능이고, 따라서
앞의 차집합도 국소 구성가능이다.

$G^\circ$가 닫혀 있고 $G\longrightarrow S$가 $G^\sigma$의 점들에서
보편적으로 열려 있다고 가정하자. $G^\sigma\longrightarrow S$가 열린
사상임을, 즉 $G^\sigma$ 안에서 $x\in G^\sigma$의 한 근방을
$y=f(x)$의 한 근방으로 보냄을 증명하자. $G^\sigma$가 국소
구성가능이므로, $y$의 모든 일반화 $y'$에 대하여 $y'$ 위에 놓이는
$G^\sigma$ 안의 $x$의 일반화 $x'$이 존재함을 보이면 된다. 밑변경을
하면 $S$가 이산 값매김환의 스펙트럼이고 $y,y'$이 각각 닫힌점과
일반점인 경우로 환원된다. $G\longrightarrow S$가 $x$에서 열려 있으므로
$y'$ 위에 놓이는 $G$ 안의 $x$의 일반화 $x_1$이 존재한다. 다시 밑변경을
할 수 있으므로, $x=g(y)$를 만족하는 $S$ 위의 $G$의 단면 $g$가 있다고
더 가정해도 된다. $k(y')$의 표수가 0이면 $y'$ 위에 놓이는 $G$ 안의
$x$의 아무 일반화 $x'$이나 택하면 된다. $G^\tau$가 열려 있으므로
$x'\in G^\tau$이고, $G_{y'}^\sigma=G_{y'}^\tau$이므로
$x'\in G^\sigma$이다. $k(y')$의 표수가 $p>0$이면
\[
\nu(g(y'))=p^hm\qquad\text{이며 }(m,p)=1
\]
로 놓는다.

\src{227}
$ap^h+bm=1$을 만족하는 두 정수 $a,b$를 잡고
\[
g_1=g^{ap^h},\qquad g_2=g^{bm},\qquad\text{따라서}\qquad g=g_1g_2
\]
로 놓자. 그러면 정의에 따라 $g_1(y')\in G^\sigma$이고
$g_2(y')\in G^\rho$이다. $G^\circ$가 닫혀 있으므로
$g_1(y)\in G^\sigma$이고 $g_2(y)\in G^\rho$이다. 그런데
\[
g(y)=g_1(y)g_2(y)\in G^\sigma
\]
이므로 $g_2(y)\in G^\sigma$이기도 하다. 따라서
\[
g_2(y)\in G_y^\sigma\cap G_y^\rho=G_y^\circ.
\]
가정과 $S$가 이산 값매김환의 스펙트럼이라는 사실에서
$G-(G_y-G_y^\circ)$는 $G$의 열린집합이고, 그 위에서
$G\longrightarrow S$가 열린 사상을 유도함을 얻는다. 따라서
$G_y^\circ$의 모든 점을 지나는 ``준단면''이 존재한다. 그러므로 필요하면
밑공간 $S$를 다시 확대하여, $g'_2(y)=g_2(y)$를 만족하는 $S$ 위의
$G^\circ$의 단면 $g'_2$가 있다고 가정해도 된다. $g'=g_1g'_2$로 놓자.
그러면 정의에 따라 $g'(y)=g(y)=x$이고
\[
g'(y')=g_1(y')g'_2(y')\in G^\sigma
\]
이다. 따라서 $g'(y')=x'$\footnote{인쇄본에는 여기에 ``$g(y')=x'$''라고
되어 있으나, 앞 문장에서 구성하고 지금 일반 섬유를 계산하는 단면은
$g'$이다.}은 $y'$ 위에 놓이는 $G^\sigma$ 안의 $x$의 일반화이다. 이로써
(iii)이 증명된다.

끝으로 (iv)의 마지막 명제를 증명하자. $x\in G^\rho$이면 $x$의 모든
일반화 $x'$이 $G^\rho$에 속함을 보이면 된다. 적당한 $h$에 대하여
$\varphi_{p^h}$로 상을 취하면\footnote{인쇄본에는 $\varphi_{n^h}$라고
되어 있다. $x\in G^\rho$이면 p.~223의 정의에 따라 표수 지수 $p$의 한
거듭제곱이 $x$를 $G^\circ$로 보낸다.} $x\in G^\circ$이라고 더 가정해도
된다. 이제 표수와 서로소인 모든 정수 $n$에 대하여 $x$는
$\varphi_n$의 상에 속한다(표수가 $n$과 서로소인 체 위의 연결 유한형
군에서 $n$제곱 사상은 전사이기 때문이다). $\varphi_n$이 열려 있으므로
$x'$도 $\varphi_n$의 상에 속한다. 더 정확히 말해, $G_y^\circ$를 포함하는
$G^\tau$의 준콤팩트 열린집합 $U$를 잡자. 그러면 $p$와 서로소인 모든
$n$에 대하여 $x'\in\varphi_n(U)$이다. $z\in U$에 대한 $\nu(z)$의
$p$와 서로소인 인수들의 공배수를 $n$으로 택하면 $x'\in G^\rho$를
얻는다.

정리 (1.1), (iii)는 다음과 같이 보충된다.

\result{따름정리}{1.5} $n>1$을 $n$제곱 준동형
$\varphi_n:G\longrightarrow G$가 보편적으로 열린(예를 들어 에탈인)
정수라 하고
\[
G^{(n)}=\varphi_n^{-1}(G^\circ)
\]
로 놓자. 모든 기하적 섬유 $G_{\bar s}$에서 부분집합

\src{228}
$\bigcup_{h\geq0}|{}_{n^h}G_{\bar s}|$가
$\bigcup_{h\geq0}G_{\bar s}^{(n^h)}$에서 조밀하다고 가정하자. 그러면
$G\longrightarrow S$는 $G^{(n^h)}$들의 점들에서 보편적으로 열려 있다.
특히 $G_s^\circ$들이 가법 성분을 포함하지 않고, 모든 정수 $n>1$에
대하여 $\varphi_n$이 잉여표수가 $n$과 서로소인 점들에서 보편적으로
열려 있다고 가정하자\footnote{인쇄본에는 ``$p$와 서로소''라고 되어 있다.
$n$으로 양화된 이 명제에서 에탈성 판정법과 핵 ${}_nG$를 이용하는 논증은
잉여표수가 $n$과 서로소일 것을 요구한다.}. 그러면
$G\longrightarrow S$는 $G^\sigma$의 점들에서 보편적으로 열려 있다.
따라서 (1.1), (iii)에 의해, $G^\circ$가 닫혀 있으면
$G^\sigma\longrightarrow S$는 보편적으로 열린다. 이 경우들에는
$s\longmapsto\dim G_s$가 $S$ 위의 국소 상수 함수이다.

실제로 가정에 따라 $\varphi_n$의 핵 ${}_nG$는 $S$ 위에서 보편적으로
열려 있다. 따라서 $G\longrightarrow S$는 ${}_nG$의 점들에서 보편적으로
열려 있고, $n$을 $n^h$으로 바꾸면 ${}_{n^h}G$의 점들에서도 보편적으로
열려 있다. 조밀성 가정은 섬유마다 위수가 $n$의 거듭제곱인 점들이
$G^{(n^h)}$들의 합집합에서 조밀하다는 것을 정확히 뜻한다. 특별한 명제의
경우, 가법 성분이 없다는 사실은 잉여체에서 가역인 정수 $n$들에 대하여
이 조밀성을 함의한다. 이로부터 $G\longrightarrow S$가 이 합집합의 모든
점에서, 특히 항등원 단면을 따라 보편적으로 열려 있음을 쉽게 얻으며,
따라서 $s\longmapsto\dim G_s$가 국소 상수 함수임을 쉽게 얻는다.

\result{비고}{1.6} 사상 $X\longrightarrow S$가 모든 열린집합을
열린집합으로 보내고, 임의의 밑변경 $S'\longrightarrow S$ 뒤에도 이 훌륭한
성질을 유지할 때 이 사상을 보편적으로 열려 있다고 한다는 것을 상기하자. 이는
또한 ($S$가 국소 뇌터이고 $X\longrightarrow S$가 국소 유한형일 때)
$X$의 모든 기약 성분이 $S$를 지배하고, 이 성질이 임의의 밑변경
$S'\longrightarrow S$ 뒤에도 보존됨을 뜻한다. 더구나 두 진술 모두에서
(위의 유한성 가정 아래) $S'$이 이산 값매김환의 스펙트럼인 밑변경
$S'\longrightarrow S$만으로 검사해도 충분하다(원한다면 이 환을 완비이고
잉여체가 대수적으로 닫힌 것으로 취할 수 있다~...). 이 정의는 $X$의 부분집합
$Z$에도 자명하게 확장된다(예를 들면 $G$의 부분집합 $G^\sigma$). 심지어
기하적 섬유가 연결인 매끄러운 사영 사상 $X\longrightarrow S$에서 시작해도
(예컨대 $S=\operatorname{Spec}\mathbf C[j]$ 위의 타원곡선 모듈러 족의
섬유제곱), $\operatorname{Pic}_{X/S}$가 $S$ 위에서 보편적으로 열려 있지
않는 것은 지극히 정상적인 현상이다. 즉
$\operatorname{Pic}_{X/S}$의 어떤 기약 성분들이 $S$의 한 점 위에만 완전히
놓일 수 있다. 이는 $X/S$의 섬유들의 네론--세베리 군의 계수(階數)가
위쪽으로 도약할 수 있다는 사실(``복소곱셈'' 현상)과 관련된다. 반면 (1.5)는
좋은 경우 $\operatorname{Pic}_{X/S}^\sigma$가(더 나아가 흔히
$\operatorname{Pic}_{X/S}^\tau$까지도 그런 듯하다) $S$ 위에서
보편적으로 열려 있음을 보장한다.

\src{229}
마지막으로 예외적으로 $G^\circ$가 열린집합이 되어 주는 유용한 경우가
다음과 같다.

\result{따름정리}{1.7} $G\longrightarrow S$가 $G^\circ$의 점들에서
보편적으로 열려 있고(cf. (1.5)), 섬유 $G_s$가 분리 가능하여 $k(s)$ 위에서
매끄럽다고 가정하자(카르티에의 한 결과에 따라 뒤의 조건은 잉여 표수가 0이면
자동으로 충족된다). 그러면 $G^\circ$는 $G$에서 열린집합이다. 더 나아가
$S$가 축소이면 $G^\circ$는 $S$ 위에서 매끄럽다(특히 평탄하다).

첫 번째 진술은 다음과 같이 더 정밀하게 할 수 있다. 주어진 $s\in S$에 대하여
$G\longrightarrow S$가 $G_s^\circ$의 점들에서 보편적으로 열려 있고
$G_s^\circ$가 $k(s)$ 위에서 분리 가능하면, $G^\circ$는
$G_s^\circ$의 근방이다(더구나 $s$에서 한 가정은 이웃한 점들에서도 계속
성립한다). 이 진술은 어떤 군 구조에도 의존하지 않으며 [2], IV, 제7절에서
증명될 것이다. (1.7)의 마지막 진술도 어떤 군 구조나 연결 성분에 관한
문제와 무관하며, [2], IV, 제5절에 주어진 평탄성 판정법의 특수한 경우이다.
그 판정법은 더 일반적으로 다음을 함의한다.

\result{따름정리}{1.8} $U$를 섬유 $G_s$의 열린 부분집합이라 하고,
$G\longrightarrow S$가 $U$의 점들에서 보편적으로 열려 있다고 하자(cf.
(1.5)). $G_s$가 $k(s)$ 위에서 분리 가능하고 $S$가 $s$에서 축소이면,
$G$는 $U$의 점들에서 $S$ 위에 평탄하다(따라서 이 경우 그 점들에서
$S$ 위에 매끄럽다).

\result{비고들}{1.9} $G$가 $S$ 위에서 분리되어 있을 때 $G^\circ$나
$G^\tau$가 언제나 $G$에서 닫혀 있는지는 알지 못한다. 그럴 가능성은 작아
보인다. 어쨌든 분리 가정을 버리면 자명한 반례들이 있다. 예컨대 아핀 직선의
원점을 가산 무한 번 중복시키면, 모든 섬유가 항등군 하나로 축약되되 한 섬유만
$\mathbf Z$인 아핀 직선 위의 군 스킴 $G$를 얻는다. 이 군 스킴은 두 직선이
한 점에서 만나도록 퇴화하는 원뿔곡선족에 대응하는 $S$-스킴 $X$의 피카르
준스킴 안에서, 항등 단면의 폐포인 열린 부분군 스킴이다.

더구나 이산 값매김환 $V$의 스펙트럼 $S$ 위의 유한 평탄 군 스킴 $G$에서
출발하여 $G^\circ$가 항등 단면 하나로 축약되고 따라서 닫혀 있는 경우에도,
몇몇 가정을 버리면 여러 결론이 거짓이 된다. $V$의 동일 표수가 $p>0$이라고
하고, 함자들의 준동형
\[
f\longmapsto f^p-tf
\]
로 정의되는 준동형 $\mathbf G_a\longrightarrow\mathbf G_a$의 핵을 $G$라
하자. 여기서 $t$는 $V$의 균등화원이다.

\src{230}
(정의에 따라 $S$ 위의 ``가법군'' $\mathbf G_a$는 함자
$S'\longmapsto\Gamma(S',\mathcal O_{S'})$를 표현함을 상기하라.) 그러면
$S$-스킴으로서
\[
G=\operatorname{Spec}V[x]/(x^p-tx)
\]
는 한 점에서 만나는 $p$개의 ``직선''의 합집합이며, 가법형 무한소군으로
퇴화하는 군 $\mathbf Z/p\mathbf Z$를 나타낸다. 이 예에서는
$G^\circ=G^\sigma$이고(이는 $p$개 직선 가운데 하나이다), 따라서
$G$에서 열려 있지 않다. 이제 잉여 표수가 $p>0$인 혼합 표수의 경우에는
$\mathbf G_m$의 $p$제곱 사상의 핵인 군 스킴 $H=\boldsymbol\mu_p$에서
출발하자. 따라서
\[
H=\operatorname{Spec}V[x]/(x^p-1) ;
\]
이것 역시 한 점에서 만나는 $p$개의 직선의 합집합이며, 표수 0에서의 군
$\mathbf Z/p\mathbf Z$가 표수 $p$에서 곱셈형 무한소군으로 퇴화하는 모습을
나타낸다. 통상적인 유한군 $\mathbf Z/p\mathbf Z$가 정의하는 ``상수 군
스킴''을 $H'$이라 하자. 이는 $\mathbf Z/p\mathbf Z$개인 $S$의 서로소 합,
달리 말해 $\operatorname{Spec}V^{\mathbf Z/p\mathbf Z}$이다. 그러면
$G=H\times_S H'$은 표수 0에서 군 $(\mathbf Z/p\mathbf Z)^2$을 나타내고,
표수 $p$에서는 무한소군과 $\mathbf Z/p\mathbf Z$의 곱으로 퇴화한다.
여기서 $G^\rho$는 항등 단면과 특수 섬유의 합집합이므로 열려 있지 않은
반면, $G^\sigma$는 항등 단면과 일반 섬유의 합집합이므로 닫혀 있지 않다.
이는 (1.1), (iii), (iv)의 동일 표수 경우에 한 주장과 대조된다. 물론 이는
표수 $p>0$과 관련된 현상이다. 실제로 앞의 결과들은 다음을 준다.

\result{따름정리}{1.10} $S$의 모든 잉여 표수가 0이라고 가정하자. 따라서
$G^\tau=G^\sigma$이고 $G^\circ=G^\rho$이다. 그러면
$G^\tau=G^\sigma$는 열려 있다. $G$가 $S$ 위에서 분리되어 있고
$G_s^\circ$들이 고유이면 이는 더 나아가 열리고 닫혀 있다. 같은 가정 아래
$G^\circ=G^\rho$는 $S$ 위에서 고유이므로 $G$에서 닫혀 있다. 마지막으로
모든 정수 $n>1$에 대하여 $G$의 $n$제곱 준동형이 보편적으로 열려 있다고
가정하면 $G^\circ$는 열려 있다. 더 나아가 $G_s^\circ$들이 가법 성분을
갖지 않으면(예컨대 위와 같이 $G_s^\circ$들이 고유이면),
$G^\tau\longrightarrow S$는 보편적으로 열려 있으며, $S$가 축소이면 더
나아가 매끄럽다.

마지막으로 다음의 쉬운 결과를 언급하자.

\result{명제}{1.11} $S$의 열린 부분집합 $U$가 존재하여, $G$가 $S$ 위에서
매끄러운(각각 평탄한) $G$의 점들의 집합 $G'$은 $G|U$에서 유도된 열린
부분군 스킴의 바탕 집합이다. 더 나아가 $U$ 위의 $G$의 모든 단면은 $U$
위의 $G'$의 단면이다.

\src{231}
\result{따름정리}{1.12} $G$가 항등 단면의 점들에서 $S$ 위에
매끄럽다면(각각 평탄하다면), $S$ 위의 $G$의 모든 단면의 점들과
$G^\circ$의 모든 점에서도 그러하다. 더 나아가 모든 정수 $n>0$에 대하여
$n$제곱 준동형 $\varphi_n:G\longrightarrow G$가 표수가 $n$과 서로소인
점들에서 에탈이면, $G$는 $G^\sigma$의 모든 점에서 $S$ 위에
매끄럽다(각각 평탄하다).

\fsection{2}{피카르 스킴의 국소적 성질에 대한 응용}

\result{정리}{2.1}

\noindent(i) $f:X\longrightarrow S$를 고유하고 매끄러운 사상이라 하고,
$\operatorname{Pic}_{X/S}$가 존재한다고 가정하자(예를 들어 $f$가 사영
사상인 경우). 그러면 $\operatorname{Pic}_{X/S}$는 $S$ 위에서 분리되어
있고, $S$ 위에서 유한형인 $\operatorname{Pic}_{X/S}$의 모든 닫힌 부분집합
$Z$는 $S$ 위에서 고유이다.

\noindent(ii) $X$를 체 $k$ 위에서 고유이고 기하적으로 정규인 준스킴이라
하자. 그러면 $\operatorname{Pic}_{X/k}^\circ$\footnote{인쇄본에는
$\operatorname{Pic}_{X/S}^\circ$라고 되어 있으나, 체 $k$ 위의 이 명제에는
$S$가 나타나지 않는다.}는 $k$ 위에서 고유이다.

\begin{proof}
(i) 값매김 판정법([2], II, 7절)에 따라 다음을 증명하면 충분하다. $S$를
완비 이산 값매김환의 스펙트럼이라 하고, $U$를 $S$의 일반점 하나로 이루어진
열린 부분집합이라 하자. 그러면 $S$ 위의 $\operatorname{Pic}_{X/S}$의 모든
유리 단면, 즉 $U$ 위의 모든 단면은 $S$ 위의 단면으로 유일하게 연장된다.
$\operatorname{Pic}_{X/S}$의 정의를 고려하면 이는 다음 명제와 동치이다.
$V=f^{-1}(U)$ 위의 모든 가역 가군 $\mathcal L$에 대하여, $\mathcal L$을
연장하는 $X$ 위의 가역 가군이 존재하고 동형사상을 제외하고 유일하다. 그런데
$V$ 위와 $X$ 위의 가역 가군을 각각 ``카르티에'' 인자류로 기술하면 이 사실을
쉽게 얻는다. 실제로 $X$의 국소환들은 정칙이다($X$가 정칙인 $S$ 위에서
매끄럽기 때문이다). 따라서 아우슬란더의 결과에 의해 이 국소환들은
유일인수분해정역이고, 이에 따라 $X$ 위의 모든 인자\footnote{인쇄본에는
여기서 ``$S$ 위''라고 되어 있으나, 사용되는 유일인수분해성은 $X$의
국소환들의 성질이다.}는 카르티에 인자이다. 실제로 $V$ 위의 모든 인자는 그
``폐포''를 취함으로써 $X$ 위의 인자로 연장할 수 있다.
\end{proof}

\result{비고}{2.2} $f$가 평탄하고 그 섬유 $X_s$들이 국소 완전교차이며
여차원 $\leq2$에서 $k(s)$ 위로 매끄럽다고만 가정해도 위의 증명은 유효하다.
이는 [5]에서 증명된 다음 결과에 따른다. 여차원 $\leq3$에서 정칙인 뇌터
국소 완전교차환은 유일인수분해정역이다(``사무엘의 추측''). 다음 사실에
유의하자.

\src{232}
``여차원 $\leq2$''를 ``여차원 $\leq1$'', 즉 ``정규''라는 가정으로
바꾸면 이 결과는 거짓이 된다. 비특이 이차 초곡면들의 족이 이차 원뿔로
퇴화하는 예에서 이를 확인할 수 있다.

(ii) 차우 보조정리를 사용하면 $X$가 사영인 경우로 환원된다. 따라서
$X\subset\mathbf P_k^n$이라 할 수 있고, $X$가 연결이라고 가정해도 된다.
$\dim X=1$이면 $X$는 $k$ 위에서 매끄러우므로 (i)을 적용한다.
$\dim X\geq2$이면 알려진 ``동치 판정법''의 한 변형을 사용할 수 있다. 이에
따르면 $\mathbf P_k^n$의 적절한 선형 부분공간들과 $X$의 교차로 얻어지는,
$X$ 위의 매끄러운 곡선 $Y_i$들이 유한 개 존재하여
\[
\operatorname{Pic}_{X/k}^\tau\longrightarrow
\prod_i\operatorname{Pic}_{Y_i/k}^\tau
\]
가 단사사상이 되고, 따라서 특히 항등원 연결 성분들 사이에서도 단사사상을
유도한다. 우변의 항등원 연결 성분은 앞의 결과에 의해 $k$ 위에서 고유이다.
또한 지금의 사상은 군 스킴의 준동형이므로 반드시 닫힌 몰입 사상이다. 따라서
$\operatorname{Pic}_{X/k}^\circ$도 $k$ 위에서 고유이다. 대수적으로 닫힌 체
위의 가환 대수군의 구조(슈발레--보렐에 의한 것)를 사용하면 까다로운 동치
판정법에 의존하지 않을 수 있다. 원점을 뺀 아핀 직선\footnote{인쇄본에는
``아핀 직선에서''라고 되어 있다. 이에 연결된 303쪽의 정오표는 ``원점을 뺀
아핀 직선''으로 고치도록 한다.}에서 $\operatorname{Pic}_{X/k}^\tau$로 가는
모든 사상이 상수임을 증명하는 문제로 환원된다. 이는
$X[t,t^{-1}]$ 위의 모든 가역 가군\footnote{인쇄본에는 $X[t]$라고 되어
있다. 이에 연결된 303쪽의 정오표는 $X[t,t^{-1}]$로 고치도록 한다.}이 $X$
위의 한 가역 가군에서 온다는 것과 동치이다. 이 사실은 잘 알려져 있고 매우
초보적이며, $X$가 $k$ 위에서 고유라는 사실조차 사용하지 않는다($X$가
정규라는 가정으로부터 곧바로 $X$가 정칙인 경우로 환원할 수 있다).

\result{따름정리}{2.3} $f:X\longrightarrow S$를 고유한 정규 사상(즉
평탄하고 기하적으로 정규인 섬유들을 갖는 사상)이라 하고,
$\operatorname{Pic}_{X/S}$가 존재한다고 가정하자. 그러면
$\operatorname{Pic}_{X/S}^\circ$는 $S$ 위에서 고유이고, 따라서
$\operatorname{Pic}_{X/S}$ 안에서 닫혀 있다. 더 나아가
$\operatorname{Pic}_{X/S}^\tau$와 $\operatorname{Pic}_{X/S}^\rho$도
닫혀 있고, ``동표수''인 경우에는 $\operatorname{Pic}_{X/S}^\sigma$도
닫혀 있다.

(1.1)과 (2.1), (ii)를 적용하면 된다.

\result{따름정리}{2.4} $f:X\longrightarrow S$를 고유하고 매끄러운
사상이라 하고, $\operatorname{Pic}_{X/S}$가 존재하며 $S$ 위에서 유한형인
스킴 $P^{(i)}$들의 서로소 합이라고 가정하자(cf. [3], V, 4.1). 그러면 각
$P^{(i)}$는 $S$ 위에서 고유이다.

이는 (2.1), (i)에 따른다. (2.2)에서 지적했듯이 $f$의 섬유들에 덜 제한적인
가정을 두어 이 결과를 일반화할 수 있지만, $f$가 정규라고만 가정하면 거짓이
된다. 이 경우에는 $\operatorname{Pic}_{X/S}^\tau$가 $S$ 위에서 유한형이라고
가정하더라도 여전히 $S$ 위에서 고유인지 나는 알지 못한다.

\src{233}
\result{정리}{2.5} $f:X\longrightarrow S$를 고유 평탄 사상이라 하고,
$\operatorname{Pic}_{X/S}$가 존재한다고 하자. 모든 정수 $n$에 대하여 이
군 준스킴의 $n$제곱 준동형을 $\varphi_n$이라 하자. 그러면 $\varphi_n$은
잉여 표수가 $n$과 서로소인 모든 점
$x\in\operatorname{Pic}_{X/S}$\footnote{인쇄본에는 ``$x\in X$''라고
되어 있으나, $\varphi_n$은 피카르 준스킴의 거듭제곱 자기준동형이다.}에서
에탈이다.

에탈 사상의 무한소적 특징짓기에 따르면 이 주장은 다음 명제와 동치이다.

\result{보조정리}{2.6} $S$를 국소 아르틴 환 $A$의 스펙트럼이라 하고, 그
극대 아이디얼 $\mathfrak m$이 $\mathfrak m^{\nu+1}=0$을 만족한다고 하자.
다음과 같이 놓자.
\[
A_{\nu-1}=A/\mathfrak m^\nu,
\qquad X_{\nu-1}=X\otimes_A A_{\nu-1}.
\]
$\mathcal L$을 $X$ 위의 가역 가군이라 하고, $\mathcal L'_{\nu-1}$을
$X_{\nu-1}$ 위의 가역 가군으로서 그 $n$차 텐서 거듭제곱이
\[
\mathcal L_{\nu-1}=\mathcal L\otimes_A A_{\nu-1}
\]
과 동형인 것이라 하자. 그러면 $n$이 $k=k(A)$의 잉여 표수와 서로소일 때,
$n$차 텐서 거듭제곱이 $\mathcal L$과 동형인 $X$ 위의 가역 가군
$\mathcal L'$이 존재한다.

보조정리를 증명하자. 다음과 같이 놓는다.
\[
V=\mathfrak m^\nu=\mathfrak m^\nu/\mathfrak m^{\nu+1}.
\]
이는 $k=k(A)$ 위의 벡터공간이다. 우선 $\mathcal L'_{\nu-1}$을 $X$ 위의
어떤 가역 가군 $\mathcal L'$으로 연장해 보자. 그러한 연장의 장애물은
\[
H^2(X_0,\mathcal O_{X_0})\otimes_k V
\]
에 속한다. 그런데 $\mathcal L'_{\nu-1}$에 둔 가정과
$\mathcal L_{\nu-1}$이 연장될 수 있다는 사실에 따라, 이 장애물에 $n$을
곱하면 영이 된다. $n$이 표수와 서로소이므로 장애물 자체가 영이다. 한편
선택한 연장의 불확정성은
\[
H^1(X_0,\mathcal O_{X_0})\otimes_k V
\]
에 속하고, $\mathcal L'^{\otimes n}$과 $\mathcal L$ 사이의 차이 $\xi$도
같은 가군에 속한다. 이 차이가 영이 되도록 $\mathcal L'$을 고치려면, 그
가군 안에서 $n\eta=\xi$를 만족하는 $\eta$를 찾으면 된다. 이는 역시 $n$이
표수와 서로소이므로 가능하다.

\result{따름정리}{} (2.5)의 조건 아래에서, 섬유 $X_s$들의 피카르 스킴이
가법 성분을 갖지 않는다고 더 가정하자(예를 들어 $X_s$들이 기하적으로
정규인 경우, cf. (2.1), (ii)). 그러면
$\operatorname{Pic}_{X/S}\longrightarrow S$는
$\operatorname{Pic}_{X/S}^\circ$의 점들에서 보편적으로 열려 있다.
$\operatorname{Pic}_{X/S}^\circ$가 닫혀 있으면(예를 들어 $X_s$들이
기하적으로 정규인 경우, cf. (2.3)),
$\operatorname{Pic}_{X/S}^\sigma$도 $S$ 위에서 보편적으로 열려 있다.
마지막으로 동표수인 경우
$\operatorname{Pic}_{X/S}^\rho\longrightarrow S$는 보편적으로 열려 있다.

(1.5)와 (1.1)을 적용하면 된다.

\src{234}
\result{따름정리}{2.7} $f:X\longrightarrow S$\footnote{인쇄본에는
``$f:X\longrightarrow Y$''라고 되어 있으나, 뒤따르는 함수와 피카르
준스킴은 모두 $S$를 밑으로 한다.}를 고유 평탄 사상이라 하고,
$\operatorname{Pic}_{X/S}$가 존재한다고 하자. 그러면 $S$ 위의 함수
\[
s\longmapsto\dim\operatorname{Pic}_{X_s/k(s)}
\]
는 상반연속이다(즉 위쪽으로 도약할 수는 있지만 아래쪽으로 도약할 수는 없다).
또한 $\operatorname{Pic}_{X_s/k(s)}$들이 가법 성분을 갖지 않으면 이 함수는
연속이기까지 하다(즉 국소 상수이다).

첫 번째 주장은 국소 뇌터 밑 위에서 국소 유한형인 모든 군 준스킴에 대하여
자명하거나 거의 자명하다. 항등 단면을 따라 살펴보는 것으로 충분하기
때문이다. 두 번째 주장은 (2.5)에 따른다.

\result{비고}{2.8} $s,s'\in S$이고 $s$가 $s'$의 특수화라고 하자. 그러면
(2.7)은 $X_{s'}$의 피카르 스킴과 그 ``특수화'' $X_s$의 피카르 스킴의
차원 사이의 부등식(각각 등식)과 동치이다. 더구나 세르는 피카르 스킴이
구성되기 전부터, 매끄러운 사상 $f:X\longrightarrow S$의 경우 $X_s$들의
피카르 스킴 차원의 불변성이 기본군의 특수화 이론([4], X)과 피카르 다양체
위의 유한 위수 점들과 아벨화한 기본군 사이의 고전적인 쿠머 관계([4], XI)의
형식적 귀결임을 지적했다. 일반적으로
$\operatorname{Pic}_{X_s/k(s)}^\circ$의 아벨 부분, 곱셈적 부분, 가법적
부분의 차원을 각각 $\alpha,\mu,\lambda$라 하고,
$\alpha',\mu',\lambda'$도 같은 방식으로 정의하자. 알려진 관계들은 다음
부등식으로 표현된다.
\[
\tag{*}\alpha+\mu+\lambda\geq\alpha'+\mu'+\lambda'
\]
(이는 $\operatorname{Pic}_{X/S}$가 존재하기만 하면 성립하고, 따라서 아마
항상 성립할 것이다). $\lambda=\lambda'=0$이면 이 부등식은 같은 존재 가정
아래에서 성립하는 등식
\[
\alpha+\mu=\alpha'+\mu'
\]
으로 바뀐다. 한편 다음 부등식도 있다.
\[
\tag{**}2\alpha+\mu\leq2\alpha'+\mu'
\]

\src{235}
$X_s$들이 분리가능이면 세르의 논증에 따라 이 부등식이 성립한다
($\operatorname{Pic}_{X/S}^\tau$의 존재조차 가정할 필요가 없다). 또는
${}_n\operatorname{Pic}_{X/S}$, 즉 피카르 함자 안에서 $\varphi_n$의 핵들이
$n$이 가역인 $S$의 열린 부분집합 위에서 분리되어 있는 경우에도 성립한다.
실제로 (2.5)에 따라 이 핵들은 그 열린 부분집합 위에서 에탈이다. (*)가 항상
등식이거나, 적어도 $X_s$들이 분리가능일 때에는 등식이라고 추측하고 싶다.
또한 다음 부등식들이
\[
\tag{***}\alpha\leq\alpha',\qquad\lambda\geq\lambda'
\]
$S$가 국소 뇌터일 때, $S$ 위에서 국소 유한형이고 섬유의 차원이 일정한 군
준스킴이 주어질 때마다 성립할 것으로 추측하고 싶다(이 방향의 긍정적인
결과는 (1.3)을 보라).

\result{비고}{2.9} 알려진 모든 경우에
$\operatorname{Pic}_{X/S}^\tau$는 $S$ 위에서 보편적으로 열려 있다. 그러나
$f:X\longrightarrow S$가 매끄럽더라도 이를 근거로 지나친 기대를 품어서는
안 될 듯하다. 어쨌든 멈퍼드는 $\operatorname{Pic}_{X/S}^\tau$가 $S$
위에서 평탄하지 않은 예를 구성했다. 이 예에서 $S$는 축소가 아니며, 실제로
아르틴 환의 스펙트럼이고, 이구사 곡면을 무한소적으로 변형하여 얻어진다.
멈퍼드가 고찰한 점은 $\operatorname{Pic}_{X/S}^\rho$에 속한다. 한편
$f:X\longrightarrow S$가 매끄러울 때 $\operatorname{Pic}_{X/S}^\tau$가
$\operatorname{Pic}_{X/S}^\sigma$의 점들에서 $S$ 위에 평탄할 가능성은
남아 있다. 그러나 강연자는 $\operatorname{Pic}_{X/S}^\circ$의 점들로
제한하고 $S$가 이산 값매김환의 스펙트럼이라고 가정하더라도 이것이 항상
성립할지는 의심한다. 이 문제는 유한 자기동형군의 작용 아래 아벨 스킴의
고정점을 연구하는 문제와도 관련되는데, 이 상황에 대해서는 예들이 부족한
듯하다.

$f:X\longrightarrow S$가 매끄럽고 사영이라고 제한하더라도, 이 절에서 말한
$\operatorname{Pic}_{X/S}$의 국소적 정칙성 결과들과 (2.8)에서 언급한
추측들은 $f$의 섬유의 성질에 더 특별한 가정을 두지 않고 이 문제에 관해 말할
수 있는 것을 거의 다 포괄하는 듯하다. 그러나
$\operatorname{Pic}_{X/S}$의 기하적 섬유들이 축소이고 가법 성분을 갖지
않으면, (1.8)과 (2.5)에 따라 $S$가 축소일 때
$\operatorname{Pic}_{X/S}$는 $\operatorname{Pic}_{X/S}^\sigma$의 점들에서
$S$ 위에 매끄럽다는 사실을 상기하자. $f:X\longrightarrow S$가 정규이면
$S$에 아무 가정도 두지 않고 같은 결과가 성립하며, 이는 (3.5)에서 보게 될
것이다. 이와 관련하여 다음을 언급하자.

\result{명제}{2.10}

\noindent(i) $\operatorname{Pic}_{X/S}$가 항등 단면의 점들에서 $S$ 위에
매끄럽다면(각각 평탄하다면), $\operatorname{Pic}_{X/S}^\sigma$의 모든
점과 $S$ 위의 $\operatorname{Pic}_{X/S}$의 모든 단면의 점들에서도 그러하다.

\src{236}
\noindent(ii) $s\in S$가
\[
H^2(X_s,\mathcal O_{X_s})=0
\]
을 만족한다고 하자. 그러면 $s$의 열린 근방 $U$가 존재하여
$\operatorname{Pic}_{X/S}|U$는 $U$ 위에서 매끄럽다.

\noindent(iii) $X$를 체 $k$ 위의 고유 스킴이라 하자. 그러면
\[
\dim\operatorname{Pic}_{X/k}\leq\dim H^1(X,\mathcal O_X)
\]
이고, 등호가 성립할 필요충분조건은 $\operatorname{Pic}_{X/k}$가 $k$ 위에서
매끄러운 것이다. $k$의 표수가 0이면 이는 항상 성립한다.

(i)은 (2.5)와 (1.12)에 따르고, (ii)는 매끄러운 사상의 무한소 판정법과 잘
알려진 장애물 계산에 따른다. 실제로 ``반연속성 정리''에 따라 $s$에서의
가정은 이웃한 점들에서도 계속 성립한다. 마지막으로 (iii)은
$H^1(X,\mathcal O_X)$가 $\operatorname{Pic}_{X/k}$의 항등원에서의 자리스키
접공간과 동형이라는 사실에 따른다. 마지막 주장은 표수 0의 ``형식군''은
$k$ 위에서 형식적으로 매끄럽다는 카르티에 정리의 특수한 경우이다.

\fsection{3}{$\operatorname{Pic}_{X/S}$의 표준 아벨 부분스킴과 알바네세
스킴}

\result{명제}{3.1} $k$를 체라 하고, $G$를 $k$ 위의 유한형 군 스킴으로서
가환이고 ``가법 성분이 없는'' 것이라 하자. 그러면
$G_{\mathrm{réd}}^\circ$는 $k$ 위에서 분리가능하며, 따라서 $k$ 위의
매끄러운 부분군 스킴이다.

$k$가 완전체일 때에는, 특히 $\bar k$가 $k$의 대수적 폐포일 때
$G_{\bar k}$에 대해서는 이 주장이 자명하므로,
$(G_{\bar k}^\circ)_{\mathrm{réd}}$가 $G$의 한 부분스킴에서 유래함을 보이면
충분하다. 그런데 $G_{\bar k}$에 대한 가정(가법 성분이 없음)\footnote{인쇄본에는
``Or l'hypothèse ...''라고 되어 있으나, 이에 연결된 정오표 303쪽은
``Or, de l'hypothèse ...''로 고치도록 규정한다.}에서 쉽게 다음이 따른다.
어떤 정수 $m$이 존재하여 $(G_{\bar k}^\circ)_{\mathrm{réd}}$는
$G_{\bar k}$에서의 $m$제곱 준동형의 ``스킴론적'' 상이다. 이 준동형은
$G$에 대한 같은 준동형에서 유래하므로, 후자의 스킴론적 상이 요구한
부분스킴이다.

\result{따름정리}{3.2} $X$를 $k$ 위에서 기하적으로 정규이고 고유인 스킴이라 하고
$\operatorname{Pic}_{X/k}$가 존재한다고 하자. 그러면
$\operatorname{Pic}_{X/k}$의 아벨 부분스킴 $A$로서 바탕 집합이
$\operatorname{Pic}_{X/k}^\circ$인 것이 존재한다.

실제로 (2.1), (ii)에 따라 $\operatorname{Pic}_{X/k}^\circ$는 $k$ 위에서
고유이므로 (3.1)의 조건을 만족한다.

\src{237}
따라서 앞의 결과는 어떤 경우에는 고전적인 ``피카르 다양체''(현대 이론에서는
$(\operatorname{Pic}_{X/k})_{\mathrm{réd}}^\circ$로 표기할 대상)가
체 $k$가 완전체라고 가정하지 않아도 ``$k$ 위에서 정의된다''는 것을 보인다.

이제 $f:X\longrightarrow S$를 고유하고 평탄한 상대 스킴이라 하고, 간단히
하기 위해
\[
\mathcal O_S\xrightarrow{\sim}f_*(\mathcal O_X)
\]
라고 하자. 또한 $\operatorname{Pic}_{X/S}$가 존재하고
$\operatorname{Pic}_{X/S}^\circ$가 $S$ 위에서 고유라고 하자. (3.3), (ii)를
위해서는 $\operatorname{Pic}_{X/S}^\circ$를 포함하고 $S$ 위에서
준사영인 $\operatorname{Pic}_{X/S}$의 한 열린부분이 존재한다고 더
가정한다. 앞에서 보았듯이 $f$가 사영이고 그 기하 섬유들이 분리가능하며
기약이면 이 조건이 성립한다. $S$ 위의 아벨 스킴이란 $S$ 위의 군
스킴으로서 $S$ 위에서 고유하고 매끄러우며 기하 섬유들이 연결인 것을
말함을 상기하자. 우리는 $\operatorname{Pic}_{X/S}$의 부분군 스킴 $A$로서
아벨 스킴이고 바탕 집합이 $\operatorname{Pic}_{X/S}^\circ$인 것이
존재하는지를 살펴보고자 한다. $S$가 한 체의 스펙트럼이면 그러한 것이
항상 존재함을 방금 보았다. 여기서 생각하는 일반적인 경우에는 다음을
말할 수 있다.

\result{정리}{3.3} 앞의 조건 아래에서 다음이 성립한다.

\noindent(i) $\operatorname{Pic}_{X/S}$의 아벨 부분스킴으로서 바탕 집합이
$\operatorname{Pic}_{X/S}^\circ$인 것이 존재하면, 그것은 유일하다. 따라서
그 형성은 밑변경과 양립한다.

\noindent(ii) 그러한 아벨 부분스킴이 존재하기 위해서는, $S'$가 아르틴
국소 스킴인 모든 밑변경 $S'\longrightarrow S$ 뒤에 그러한 것이 존재하면
충분하다. $S$가 한 국소환의 스펙트럼이면
$A_n=A/\mathfrak m^{n+1}$로 놓았을 때
$S'=\operatorname{Spec}(A_n)$인 경우들만 시험해도 충분하다. $S$가
축소이면, $S'$가 이산 값매김환의 스펙트럼인 밑변경
$S'\longrightarrow S$들만 시험해도 마찬가지로 충분하다(원한다면 그 환은
완비이고 그 잉여체는 대수적으로 닫혀 있다고 해도 된다).

\noindent(iii) $A$가 존재한다고 가정하고,
$B=\operatorname{Alb}^0(X/S)$를 그 쌍대 아벨 스킴, 즉
$B=\operatorname{Pic}_{A/S}^\circ$ [7]이라 하자. 그러면 $B$에 대한
주동차 공간 $P=\operatorname{Alb}^1(X/S)$와 $S$-사상
$X\longrightarrow P$를 표준적으로 구성할 수 있으며, 이 사상은 $X$에서
준아벨 스킴들(즉 아벨 스킴에 대한 주동차 공간들)로 가는 $S$-사상들에
대하여 보편적이다. $\operatorname{Alb}^0(X/S)$,
$\operatorname{Alb}^1(X/S)$ 및
$X\longrightarrow\operatorname{Alb}^1(X/S)$의 형성은 밑변경과 가환한다.

증명을 개략적으로 설명하자.

(i)은 가환 군 스킴들의 아벨 부분스킴에 관한 일반적인 강직성 성질이다.
그러한 두 부분스킴이 한 점 $s\in S$에서 집합론적으로 일치하면, $s$가
속한 연결 성분

\src{238}
전체 위에서 서로 일치한다 [7]. (이 결과는 Chow의 고전적인 한 정리를
일반화한다.)

(ii) 힐베르트 스킴을 사용하면, $S$ 위의 각 $S'$에
$(\operatorname{Pic}_{X/S})\times_S S'$의 표준 아벨 부분스킴들의 집합
(원소가 하나 또는 0개인 집합)을 대응시키는 함자가 $S$ 위의 유한형 스킴
$T$로 표현됨을 알 수 있다. (i)에 따라 $T\longrightarrow S$는
단사사상이고, (3.2)에 따라 전사이다. $\operatorname{Pic}_{X/S}$의 표준
아벨 부분스킴이 존재한다는 것은 $T$가 $S$ 위의 절단을 갖는다는 것, 다시
말해 $T\longrightarrow S$가 동형사상이라는 뜻이다. 이는
$T\longrightarrow S$가 에탈이라는 것과도 동치이고, $S$가 축소인
경우에는 $T\longrightarrow S$가 고유하다는 것과도 동치이다. 이로부터
(ii)가 곧 따른다.

(iii) $\operatorname{Pic}_{X/S}$의 정의만을 사용하면, $S$ 위의 임의의
아벨 스킴 $C$에 대하여 $X$에서 $C$에 대한 한 주동차 공간으로 가는
$S$-사상을 주는 것은 군 준동형
$C'\longrightarrow\operatorname{Pic}_{X/S}$를 주는 것과 동치임을 알 수
있다. 여기서 $C'$는 $C$의 쌍대 아벨 스킴이다. 그런데
$\operatorname{Pic}_{X/S}$의 표준 아벨 부분스킴 $A$가 존재하면, 이러한
준동형들은 반드시 $A$를 통해 인수분해된다(유한 위수의 점들을 사용해서도
이를 알 수 있다). 이로부터 (iii)가 곧 따른다.

\result{비고}{3.4} $\operatorname{Pic}_{X/S}^{\circ\circ}$로
$\operatorname{Pic}_{X/S}$의 표준 아벨 부분스킴을, 그것이 존재할 때,
나타내겠다. 불행히도 이것은 항상 존재하지는 않으며, 이구사 곡면을
무한소적으로 변형시키면(1차 모듈라이 변형을 통해) 이를 알 수 있다.
그럼에도 $\operatorname{Pic}_{X/S}^{\circ\circ}$는 적어도 $S$가 축소일
때에는 존재할 가능성이 있다. 또는 (ii)에 따르면 같은 말이지만, $S$가
한 이산 값매김환의 스펙트럼일 때에는 존재할 가능성이 있다. $X_0,X_1$을
$X$의 특수섬유와 일반섬유라 하고,
\[
A_1=\operatorname{Pic}_{X_1/K}^\circ
\]
라 하자. 여기서 $K$는 값매김환 $V$의 분수체이다. Koizumi의 결과에 따라
일반섬유가 $A_1$인, $S$ 위의 본질적으로 유일한 아벨 스킴 $A$가 존재한다.
이제 $X$가 $S$ 위에서 매끄럽다고 가정하면, (2.1), (i)에서와 마찬가지로
$A_1$의 항등사상이 다음 사상으로 연장됨을 쉽게 알 수 있다.
\[
A\longrightarrow\operatorname{Pic}_{X/S}.
\]
따라서 준동형
\[
\tag{*}A_0\longrightarrow\operatorname{Pic}_{X_0/k}^{\circ\circ}
\]
을 얻는다. 이것은 핵이 유한 $p$-군인 전사 준동형임을 쉽게 보일 수 있다.
여기서 $p$는 잉여체 $k$의 표수 지수이다.

\src{239}
(여기서도 유한 위수의 점들을 사용한다.) 이제 $X/S$에 대한 다음 조건들은
서로 동치이다.

\medskip
\noindent a. 앞의 준동형 (*)는 동형사상이다(Shimura의 표현으로는
``피카르 다양체''의 형성이 ``특수화와 양립한다'').

\noindent b. $\operatorname{Pic}_{X/S}^{\circ\circ}$가 존재한다(그때 그것은
$A$와 다름없다).

\noindent c. (기록을 위해) 모든
$\operatorname{Pic}_{X_n/S}^{\circ\circ}$가 존재한다.

(*)의 핵에 관해 한 비고에 따르면, 잉여 표수가 0이면 조건 (a)가 성립한다.
그러나 이 결과는 (3.5)에서 상당히 일반화될 것이다.

물론 $\operatorname{Pic}_{X/S}$가 $\operatorname{Pic}_{X/S}^\circ$의
점들에서 $S$ 위로 매끄러우면, 후자는 $\operatorname{Pic}_{X/S}$에서
열려 있다(cf. (1.7)). 따라서 유도된 구조를 주면 그것은
$\operatorname{Pic}_{X/S}$의 아벨 부분스킴이고, 그러므로 이 경우에
존재하는 $\operatorname{Pic}_{X/S}^{\circ\circ}$와 같다. 그러나 훨씬
더 좋은 결과가 성립한다.
\result{정리}{3.5} (3.3)의 조건 아래, $s\in S$이고
$\operatorname{Pic}_{X_s/k(s)}$가 $k(s)$ 위에서 매끄럽다고 하자(또는,
동치인 조건으로,
\[
\dim\operatorname{Pic}_{X_s/k(s)}=\dim H^1(X_s,\mathcal O_{X_s})).
\]
그러면 $s$의 열린 근방 $U$가 존재하여
$\operatorname{Pic}_{X/S}$는 $\operatorname{Pic}_{X/S}^\circ|U$의
점들에서 $S$ 위로 매끄럽다. 따라서 $\operatorname{Pic}_{X/S}^\circ|U$는
$\operatorname{Pic}_{X/S}|U$의 열린 아벨 부분스킴이다. 특히
$\operatorname{Pic}_{X|U/U}^{\circ\circ}$가 존재한다.

증명의 원리를 설명하자. 앞의 논의에 따라 $S$가 아르틴 국소환 $A$의
스펙트럼인 경우로 환원할 수 있고, $A$의 무한소 차수에 관하여 귀납한다.
따라서 $\operatorname{Pic}_{X_n/A_n}^\circ$가 $A_n$ 위에서 매끄럽다고
가정하고, $\operatorname{Pic}_{X_{n+1}/A_{n+1}}^\circ$가 $A_{n+1}$ 위에서
매끄러움을 보이는 것으로 충분하다. 이를 위해서는
\[
P_n=\operatorname{Pic}_{X_n/A_n}^\circ
\]
을 연장하는 $A_{n+1}$ 위의 아벨 스킴 $P_{n+1}$과 동시에,
$X_n\times_{A_n}P_n$ 위의 가역층 $\mathcal L_n$을 연장하는
$X_{n+1}\times_{A_{n+1}}P_{n+1}$ 위의 가역층 $\mathcal L_{n+1}$을
구성하면 충분함을 확인할 수 있다. 여기서 $\mathcal L_n$은 보편 문제의
해로서 피카르 스킴 $\operatorname{Pic}_{X_n/A_n}$을 정의할 때 나타나는
가역층이다. (주의 --- $X$가 $S$ 위의 단면을 갖는다고 가정할 수 있다~\ldots)
이 구성을 위해 다음 핵심 결과를 써야 한다. 아르틴 국소환의 몫환 위에
정의된 모든 아벨 스킴은 그 아르틴 국소환 위로 연장될 수 있다. 다시 말해,
대수적으로 닫힌 체 위의 아벨 스킴에 대한 절대적인 ``형식적 모듈라이
스킴''([3], II)은 $k$ 위의

\src{240}
비트 벡터 환 위에서 매끄럽다. 이 결과 자체는
올림 장애의 일반적인 형식적 성질과 군 연산들을 사용하여 얻는다. 이 결과를
인정하고 먼저 $P_n$을 임의의 방식으로 $P_{n+1}$로 연장한다. 그러면
$\mathcal L_n$을 올리는 데 대한 장애가 나타나며, 이 장애는
\[
H^2(X_0\times P_0,\mathcal O_{X_0\times P_0})\otimes_k V
\qquad(\text{여기서 }V=\mathfrak m^{n+1}/\mathfrak m^{n+2})
\]
에 속하고, 실제로는 그 부분공간
\[
H^1(X_0,\mathcal O_{X_0})\otimes
H^1(P_0,\mathcal O_{P_0})\otimes_k V
\]
에 속한다($\mathcal L_n$을 두 인자 $X_n$과 $P_n$에 제한하면 자명하기
때문이다). 그런데 이 마지막 공간은 바로
\[
H^1(P_0,\mathcal T_{P_0/k})\otimes V
\]
이다. 여기서 $\mathcal T_{P_0/k}$는 $P_0/k$의 접다발이다. 따라서 이는
$P_n$을 $P_{n+1}$로 올릴 때의 비유일성을 나타내는 공간이기도 하다([3],
II). 그러므로 이 올림을 (마땅히 그래야 하듯 유일한 한 가지 방식으로)
수정하여 $\mathcal L_n$을 올리는 데 대한 장애를 없앨 수 있다.

\result{따름정리}{3.6} (3.5)의 조건 아래,
$R^1f_*(\mathcal O_X)$는 $s$의 근방에서 $S$ 위의 국소 자유 가군이고,
그 구성은 밑변환과 가환한다.

실제로 이 가군은 $\operatorname{Pic}_{X/S}$의 단위원 단면을 따른
접가군이다.

\result{비고}{3.7} (3.5)에서와 같은 논증을 사용하면 다음을 보일 수 있다.
$S'$이 연접인 멱영 아이디얼층으로 정의되는 $S$의 부분스킴이고,
$\operatorname{Pic}_{X'/S'}^\circ$가 존재하며 그 섬유들이 매끄럽다고만
가정하면, $\operatorname{Pic}_{X/S}^\circ$는 반드시 존재하며 $S$ 위의
아벨 스킴이다. 이에 따라 사영성 가정이 없는 일부 경우에도 피카르 스킴을
구성할 수 있다. 예를 들어 아르틴 환 위의 아벨 스킴의 쌍대 아벨 스킴은
항상 존재한다. 한편 $X$가 $S$ 위의 아벨 스킴인 경우에 (3.6)을 적용하고,
$H^1(X_s,\mathcal O_{X_s})$ 위의 외대수로서
$H^*(X_s,\mathcal O_{X_s})$가 갖는 알려진 구조(로젠리히트--세르)를
사용하면, 모든 $i$에 대하여 $R^if_*(\mathcal O_X)$가 국소 자유이고,
더 정확히는 $R^1f_*(\mathcal O_X)$의 $i$차 외승과 동형임을 얻는다.

\result{비고}{3.8} $f:X\longrightarrow S$가 매끄러운 사영 사상이고,
$S$가 축약이며 모든 잉여 표수가 영인 경우, (3.6)의 결과는 호지 이론의
귀결로서 초월적 방법에 의해 알려져 있었다. 실제로 모든
$R^pf_*(\Omega^q_{X/S})$가 국소 자유이다. 반면 혼합 표수인 경우에는

\src{241}
세르 다양체를 이용하여 ``$R^1f_*(\mathcal O_X)$가 국소 자유이다''라는
명제의 반례를 얻는다([8], p.~50). 동표수인 경우에는 반례가 하나도 알려져
있지 않은 듯하다.

\result{따름정리}{3.8} (3.5)의 조건 아래,
$\operatorname{Pic}_{X/S}|U$는 $\operatorname{Pic}_{X/S}^\sigma|U$의
모든 점에서 $U$ 위로 매끄럽다.

(2.10)의 (i)를 적용하면 된다. 특히 (2.10)의 (iii)을 고려하면 다음을 얻는다.

\result{따름정리}{3.9} $S$의 모든 잉여 표수가 영이라고 가정하자. 그러면
$\operatorname{Pic}_{X/S}^\tau$는 $S$ 위에서 매끄럽다.

\result{비고}{3.10} 예를 들어 $\operatorname{Pic}_{X/S}^\tau$가 $S$
위에서 고유이면, $f$의 기하적 섬유들의 네론--세베리 군의 꼬임 부분군은
$S$의 각 연결 성분 위에서 일정하다는 결론을 얻는다($f$가 매끄럽고
사영이면 이는 초월적 방법으로도 명백하다). (2.5)를 직접 사용하면 더
일반적으로 다음을 보일 수 있다. $\operatorname{Pic}_{X/S}^\tau$가 $S$
위에서 고유이고(예를 들어 $f:X\longrightarrow S$가 매끄럽고 사영이면),
$q$가 $S$의 모든 잉여 표수와 다른 소수이면, $X/S$의 기하적 섬유들의
네론--세베리 군의 꼬임 부분군의 $q$-일차 성분은 $S$의 각 연결 성분
위에서 일정하다. 그러나 표수 $p>0$에서는 꼬임 부분군의 $p$-일차 성분에
대하여 더 이상 같은 결과가 성립하지 않는다. 그렇지만
\[
\operatorname{Pic}_{X_s/k(s)}^\tau/
\operatorname{Pic}_{X_s/k(s)}^{\circ\circ}=T_{X_s/k(s)}
\]
의 $k(s)$ 위의 계수가 국소 상수일 가능성은 남아 있다. $S$가 축약이면,
이것은 $\operatorname{Pic}_{X/S}^{\circ\circ}$가 존재하고
$\operatorname{Pic}_{X/S}^\tau$가 $S$ 위에서 평탄하다는 것과 동치임을
보일 수 있으며, $S$가 이산 값매김환의 스펙트럼인 경우에만 이 문제를
시험하면 충분하다. 내가 살펴본 몇 가지 예에서는 이를 확인했다. 그러나
$S$가 아르틴인 경우의 대응하는 명제는 거짓이므로((2.9)와 (3.4)의 비고
참조), 아마 지나친 기대를 품어서는 안 될 것이다.

\fsection{4}{피카르 스킴의 유한성 정리}

$f:X\longrightarrow S$를 $\operatorname{Pic}_{X/S}$가 존재하는 사영 평탄
사상이라 하자. ``힐베르트 다항식'' $Q\in\mathbf Q[t]$들을 생각하면
$\operatorname{Pic}_{X/S}$를 열린 부분스킴
$\operatorname{Pic}_{X/S}^Q$들의 서로소 합으로 분해할 수 있다. 예를 들어
$\operatorname{Pic}_{X/S}$가 $S$ 위에서 분리임을 보장하는 추가 가정을
두지 않으면, 일반적으로 이 열린 부분스킴들은 $S$ 위에서 유한형이지 않다.
$X$가 ``두 직선으로 퇴화하는 이차곡선''일 때 반례를 얻는다.

\src{242}
그렇지만 $f$가 분리가능이고 기하학적으로 기약인 섬유들을 가지면 이들이
유한형일 가능성이 있다. 이 문제는 $\operatorname{Pic}_{X/S}^\tau$가 $S$
위에서 유한형인지의 문제와 관련되며, 후자는 $X/S$의 섬유에 아무런 가정도
두지 않아도 참일지 모른다. $f:X\longrightarrow S$가 매끄러우면
$\operatorname{Pic}_{X/S}^\tau$가 $\operatorname{Pic}_{X/S}^Q$ 가운데
하나에 포함됨에 유의하자. 이는 비특이 사영 다양체 위에서 인자에 관한
``꼬임'' 동치가 수치적 동치보다 더 세밀하기 때문이다(실제로 Matsusaka에
따르면 두 동치는 ``동일하다''). 따라서 $\operatorname{Pic}_{X/S}^Q$들이
$S$ 위에서 유한형이면 $\operatorname{Pic}_{X/S}^\tau$도 그러하다. 방금
제기한 유한성 문제들은 $\operatorname{Pic}_{X/S}$의 존재조차 가정하지
않아도 명백한 의미를 갖는다. 실제로 이 문제들은 어떤 가역층들의 족이
[3]의 IV의 의미에서 ``유계''라고 말하는 것으로 표현된다. 모든 대수적으로
닫힌 체 $k$에 대하여, 차원이 $n$이고 차수가 $d$인 $\mathbf P_k^r$의 정역
부분스킴들과, 힐베르트 다항식이 $Q$인 그러한 준스킴 위의 가역층들을
생각하자(여기서 $r,n,d,Q$는 주어져 있다). 이 가역층들을
$\mathbf P_k^r$ 위의 연접 가군으로 볼 때, 그 족이 유계임, 즉 $\mathbf Z$
위의 유한형 스킴으로 매개될 수 있음을 보여야 한다.

Matsusaka의 방법 [6]을 사용하면, 상당히 기술적인 증명(``동치 판정법''을
알맞은 형태로 사용한다)을 통해 $\mathbf P_k^r$의 비특이 부분다양체들로
한정할 때에는 이 문제에 긍정적으로 답할 수 있다. 더 정확히는 다음 결과를
얻는다.
\result{정리}{4.1} $f:X\longrightarrow S$\footnote{인쇄본에는
``다음을 두자: $X\longrightarrow S$''라고 되어 있으나, 이하에서 사용하는
사상은 $f$이다.}를 기하학적으로 연결된 섬유를 갖는 매끄러운 사영 사상이라
하고, $S$는 뇌터라고 하자. $\mathcal O_X(1)$을 $S$에 관해 $X$ 위에서 매우
풍부한 가군이라 하자. $E$를 $\operatorname{Pic}_{X/S}$의 부분집합이라 하며,
이는 $X/S$의 기하학적 섬유 $\overline X_{s_i}$ 위의 가역 가군들의 족
$(\mathcal L_i)$에 대응한다고 하자. $D_i$를 $\overline X_{s_i}$ 위에서
$\mathcal L_i$를 정의하는 인자(반드시 유효일 필요는 없다)라 하고,
\[
a_n^{(i)}t^n+\cdots+a_0^{(i)}
\]
를 $\mathcal L_i$의 힐베르트 다항식이라 하며, $\xi=\xi_i$를
$\mathcal O_X(1)$을 정의하는 인자, 즉 초평면 절단이라 하자. 다음 조건들은
서로 동치이다.

\noindent a. $E$\footnote{인쇄본에는 여기서 ``$Q$''라고 되어 있다. 명제에서
도입되었으며 그 준콤팩트성이 이 족의 유계성을 나타내는 부분집합은
$E$이다.}는 준콤팩트이다. 즉 $S$ 위에서 유한형인 열린집합에 포함된다.
다시 말해 $(\mathcal L_i)$의 족은 유계이다.

\noindent b. $E$는 유한 개의 집합 $\operatorname{Pic}_{X/S}^Q$,
$Q\in\mathbf Q[t]$의 합집합에 포함된다. 즉 $\mathcal L_i$들의 힐베르트
다항식들로 이루어진 집합은 유한하다.

\noindent b'. ($X/S$의 모든 섬유의 차원이 같은 $n$인 경우.)
$\mathcal L_i$들의 힐베르트 다항식의 두 계수 $a_{n-1}^{(i)}$와
$a_{n-2}^{(i)}$는 유한한 집합 안에서만 변한다.

\src{243}

\result{따름정리}{4.2} $f:X\longrightarrow S$를 기하학적으로 연결된 섬유를
갖는 매끄러운 사영 사상이라 하자. 그러면 스킴
$\operatorname{Pic}_{X/S}^Q$와 $\operatorname{Pic}_{X/S}^\tau$
($Q\in\mathbf Q[t]$)는 $S$ 위에서 사영이다.

이 스킴들은 (4.1)에 따라 $S$ 위에서 유한형이므로 (2.4)를 적용할 수 있다.

\section*{참고문헌}
\addcontentsline{toc}{section}{참고문헌}

\begin{description}
\item[{[1]}] CHEVALLEY (Claude). -- Sur la théorie de Picard,
\emph{Amer. J. of Math.}, 제82권, 1960, 435--490쪽.

\item[{[2]}] GROTHENDIECK (A.) 및 DIEUDONNÉ (J.). -- \emph{Éléments de
géométrie algébrique}, 제I장 이하. -- Paris, Presses universitaires de France,
1960--1961 (Institut des Hautes Études Scientifiques, 4, 8, 11, \ldots).

\item[{[3]}] GROTHENDIECK (Alexander). -- Technique de descente et théorèmes
d'existence en géométrie algébrique, I--V, \emph{Séminaire Bourbaki}, 제12권,
1959/60, 제190호, 29쪽 및 제195호, 22쪽; 제13권, 1960/61, 제212호, 20쪽 및
제221호, 28쪽; 제14권, 1961/62, 제232호, 19쪽.

\item[{[4]}] GROTHENDIECK (Alexander). -- \emph{Séminaire de géométrie
algébrique}, 1960/61. -- Paris, Institut des Hautes Études Scientifiques, 1961.

\item[{[5]}] GROTHENDIECK (Alexander). -- \emph{Séminaire de géométrie
algébrique}, 1961/62, 청강자 집단이 정리함. -- Paris, Institut des Hautes
Études Scientifiques (출간 예정).

\item[{[6]}] MATSUSAKA (T.). -- The criteria for algebraic equivalence and the
torsion group, \emph{Amer. J. Math.}, 제79권, 1957, 53--66쪽.

\item[{[7]}] MUMFORD (D.) 및 TATE (J.). -- \emph{Séminaire de géométrie
algébrique}, Harvard University, 1962년 봄 학기 (출간 예정).

\item[{[8]}] SERRE (Jean-Pierre). -- Sur la topologie des variétés algébriques
en caractéristique $p$, \emph{Symposium internacional de topología algebraica}
[1956, Mexico], 24--53쪽. -- Mexico, Universidad nacional autónoma, 1958.
\end{description}

% FGA 정본 인쇄문 보존 장치: 원자적 통합 R1.
\par\smallskip
\begingroup\footnotesize
\noindent\textit{FGA-CANON-ERR-0033. 인쇄본의 독법은 C096에 계속 기록되어 있으며, FGA-CANON-DISP-0001이 수정된 현행 본문의 최소 수정을 규율한다.}\par
\noindent\textit{FGA-CANON-ERR-0028. 인쇄본의 독법은 C091에 계속 기록되어 있으며, FGA-CANON-DISP-0002가 수정된 현행 본문의 최소 수정을 규율한다.}\par
\noindent\textit{FGA-CANON-ERR-0037. 인쇄본의 독법은 C100에 계속 기록되어 있으며, FGA-CANON-DISP-0020이 수정된 현행 본문의 최소 수정을 규율한다.}\par
\noindent\textit{FGA-CANON-ERR-0034. 인쇄본의 독법은 C097에 계속 기록되어 있으며, FGA-CANON-DISP-0003이 수정된 현행 본문의 최소 수정을 규율한다.}\par
\noindent\textit{FGA-CANON-ERR-0035. 인쇄본의 독법은 C098에 계속 기록되어 있으며, FGA-CANON-DISP-0004가 수정된 현행 본문의 최소 수정을 규율한다.}\par
\noindent\textit{FGA-CANON-ERR-0025. 인쇄본의 b'' 절은 C086에 계속 기록되어 있으며, FGA-CANON-DISP-0027은 Stacks의 반례 D000571 때문에 이를 수정된 현행 본문의 동치 조건 목록에서 제외한다.}\par
\noindent\textit{FGA-CANON-ERR-0050. 가법 성분이 없으면 어떤 $n$의 거듭제곱에 의해 소멸하는 꼬임점들이 조밀하다는 인쇄본의 충분조건은 C116에 계속 기록되어 있다. FGA-CANON-DISP-0050은 섬유별 조밀성을 요구하며, 이 충분조건은 $n$이 가역인 곳에서만 유지한다.}\par
\noindent\textit{FGA-CANON-ERR-0051. ${}_n\operatorname{Pic}_{X/S}$가 $S$ 전체에서 에탈이라는 인쇄본의 주장은 C117에 계속 기록되어 있다. FGA-CANON-DISP-0051은 이 주장을 $n$이 가역인 열린집합으로 제한한다.}\par
\noindent\textit{FGA-CANON-DISP-0056부터 0058까지. 3.2에서 인쇄본의 독법은 기하학적 정규성을 요구하지 않고 $X$가 $k$ 위에서 정규이고 고유라고 가정한다(C122). 3.10에서 인쇄본의 두 고유성 가정은 $\operatorname{Pic}_{X/S}^{\tau}$가 아니라 피카르 스킴 전체에 관한 것이다(C124). C123은 2.10(i)와 3.8에서 인쇄본의 $\sigma$를 복원한다. $\circ$와 $\tau$는 필사 오류였지 인쇄본의 독법이 아니었다.}\par
\endgroup
